\documentclass[12pt,a4paper]{article}

\usepackage[utf8]{inputenc}
\usepackage[T1]{fontenc}
\usepackage[french]{babel}
\usepackage[margin=2.5cm]{geometry}
\usepackage{graphicx}
\usepackage{xcolor}
\usepackage{listings}
\usepackage{booktabs}
\usepackage{tabularx}
\usepackage{hyperref}
\usepackage{fancyhdr}
\usepackage{titlesec}
\usepackage{enumitem}
\usepackage{mdframed}
\usepackage{tcolorbox}
\usepackage{fontawesome5}
\usepackage{parskip}
\usepackage{float}

% ── Couleurs simplifiées ─────────────────────────────────────────────────────
\definecolor{darkbg}{HTML}{0F1117}
\definecolor{accent}{HTML}{3498DB}
\definecolor{gray}{HTML}{7F8C8D}
\definecolor{lightgray}{HTML}{BDC3C7}
\definecolor{codebg}{HTML}{1A1D27}
\definecolor{codetext}{HTML}{ABB2BF}
\definecolor{screenshotbg}{HTML}{F0F4F8}
\definecolor{screenshotborder}{HTML}{3498DB}

% ── Liens ─────────────────────────────────────────────────────────────────────
\hypersetup{
  colorlinks=true,
  linkcolor=black,
  urlcolor=accent,
  pdftitle={Analyseur de fichiers suspects},
  pdfauthor={Groupe Cybersécurité}
}

% ── En-tête / pied de page ────────────────────────────────────────────────────
\pagestyle{fancy}
\fancyhf{}
\rhead{\textcolor{gray}{\small Analyseur de fichiers suspects}}
\lhead{\textcolor{gray}{\small Cybersécurité défensive}}
\rfoot{\textcolor{gray}{\small \thepage}}
\lfoot{\textcolor{gray}{\small \today}}
\setlength{\headheight}{15pt}
\renewcommand{\headrulewidth}{0.4pt}
\renewcommand{\footrulewidth}{0.4pt}

% ── Style des titres ──────────────────────────────────────────────────────────
\titleformat{\section}
{\large\bfseries\color{accent}}
{\thesection.}{0.6em}{}[\titlerule]

\titleformat{\subsection}
{\normalsize\bfseries\color{accent}}
{\thesubsection.}{0.5em}{}

% ── Style de code ─────────────────────────────────────────────────────────────
\lstdefinestyle{ps1style}{
  language=sh,
  backgroundcolor=\color{codebg},
  basicstyle=\ttfamily\footnotesize\color{codetext},
  keywordstyle=\color{accent}\bfseries,
  commentstyle=\color{gray}\itshape,
  stringstyle=\color{accent},
  breaklines=true,
  frame=single,
  framesep=6pt,
  rulecolor=\color{codebg},
  xleftmargin=10pt,
  xrightmargin=10pt,
}

\lstdefinestyle{pystyle}{
  language=Python,
  backgroundcolor=\color{codebg},
  basicstyle=\ttfamily\footnotesize\color{codetext},
  keywordstyle=\color{accent}\bfseries,
  commentstyle=\color{gray}\itshape,
  stringstyle=\color{accent},
  breaklines=true,
  frame=single,
  framesep=6pt,
  rulecolor=\color{codebg},
  xleftmargin=10pt,
  xrightmargin=10pt,
}

% ── Boîtes colorées simplifiées ───────────────────────────────────────────────
\tcbuselibrary{skins, breakable}

\newtcolorbox{riskbox}[2][]{
  colback=darkbg!95!white,
  colframe=accent,
  coltext=white,
  coltitle=white,
  fonttitle=\bfseries\small,
  title=#1,
  left=8pt, right=8pt, top=6pt, bottom=6pt,
  breakable
}

% ── Boîte capture d'écran ─────────────────────────────────────────────────────
% Usage : \screenshot{nom_fichier.png}{Légende descriptive}
\newcommand{\screenshot}[2]{%
  \begin{figure}[H]
    \centering
    \begin{tcolorbox}[
      colback=screenshotbg,
      colframe=screenshotborder,
      fonttitle=\bfseries\small\color{white},
      coltitle=white,
      title={\faCamera\quad Capture d'écran — \texttt{#1}},
      left=6pt, right=6pt, top=6pt, bottom=6pt,
      width=\textwidth,
      arc=4pt
    ]
    \centering
    \vspace{0.4cm}
    \textcolor{gray}{\faImage}\\[4pt]
    {\small\color{gray}\itshape Insérer ici la capture d'écran :\\[2pt]
    \textbf{\texttt{#1}}}
    \vspace{0.4cm}
    \end{tcolorbox}
    \caption{#2}
  \end{figure}%
}

% ══════════════════════════════════════════════════════════════════════════════
\begin{document}

% ── Page de titre ─────────────────────────────────────────────────────────────
\begin{titlepage}
  \centering
  \vspace*{2cm}

  {\Huge\bfseries\color{accent} Analyseur de fichiers suspects}

  \vspace{0.5cm}
  {\large\color{gray} Rapport de cybersécurité défensive}

  \vspace{3cm}
  \begin{tcolorbox}[colback=darkbg, colframe=accent, width=0.85\textwidth,
    arc=8pt, left=20pt, right=20pt, top=16pt, bottom=16pt]
    \color{white}
    \centering
    \large
    Pipeline d'investigation : \textbf{PowerShell + Python}\\[6pt]
    \normalsize\color{lightgray}
    Inventaire des fichiers, détection des signaux faibles,\\
    empreintes SHA-256 et vérification avec VirusTotal
  \end{tcolorbox}

  \vspace{3cm}

  \begin{tabular}{rl}
    \textbf{Réalisé par} & Yassine KOOLI et Jihene GHARSALLI \\[4pt]
    \textbf{Date}        & \today \\[7pt]
  \end{tabular}

  \vfill
  {\small\color{gray}
    Tous nos droits sur ce document sont réservés. Ce dernier est destiné au seul destinataire autorisé. Sans notre autorisation écrite préalable, il ne peut être ni copié, ni reproduit, ni communiqué à des tiers.
  }
\end{titlepage}

% ── Table des matières ────────────────────────────────────────────────────────
\tableofcontents
\newpage

% ══════════════════════════════════════════════════════════════════════════════
\section{Contexte et objectifs}

\subsection{Contexte}

Lors d'une investigation en cybersécurité, l'analyste doit souvent traiter rapidement un dossier rempli de fichiers dont l'origine n'est pas certaine. Ouvrir chaque élément à la main n'est ni efficace ni fiable : on peut oublier un détail, mal comparer deux fichiers ou perdre la trace des vérifications déjà faites.

Ce rapport présente donc un petit pipeline d'analyse conçu pour rendre ce travail plus régulier. L'idée est simple : inventorier les fichiers, calculer leurs empreintes, repérer les comportements suspects les plus courants, puis produire un rapport exploitable sans devoir relire tous les résultats bruts.

\subsection{Objectifs}

\begin{itemize}[itemsep=4pt]
  \item Automatiser la collecte des métadonnées avec PowerShell sous Windows.
  \item Calculer une empreinte SHA-256 pour identifier chaque fichier de manière stable.
  \item Appliquer, côté Python, plusieurs règles simples mais complémentaires.
  \item Comparer les hashes avec une base locale et, si nécessaire, avec VirusTotal.
  \item Produire des sorties lisibles pour l'analyste : texte, JSON et HTML.
\end{itemize}

\subsection{Périmètre}

\begin{riskbox}[Environnement d'exécution]{accent}
  \color{white}
  L'outil est prévu pour un \textbf{laboratoire contrôlé}. Les scripts restent
  \textbf{non destructifs} : ils ne modifient pas les fichiers analysés et ne désactivent
  aucune protection système. Aucune donnée sensible, comme des mots de passe, tokens
  ou cookies, n'est collectée.
\end{riskbox}

% ══════════════════════════════════════════════════════════════════════════════
\section{Architecture du projet}

\subsection{Structure des fichiers}

\begin{verbatim}
python-powershell/
|-- collect.ps1                   Script PowerShell de collecte
|-- analyze.py                    Script Python d'analyse
|-- .env                          Cle API VirusTotal (non versionnee)
|-- data/
|   |-- inventory.csv             Inventaire au format tableur
|   |-- inventory.json            Inventaire JSON (entree Python)
|   |-- known_malware_hashes.json Base de donnees locale de hashes
|   `-- skipped_files.txt         Fichiers inaccessibles (journalises)
`-- results/
    |-- risk_report.txt           Rapport lisible par l'analyste
    |-- risk_report.json          Rapport structure complet
    |-- risk_report.html          Rapport interactif avec interface web
    `-- hashes.txt                Liste SHA-256 (format sha256sum)
\end{verbatim}

\screenshot{01\_structure\_projet.png}{Arborescence du projet affichée dans l'explorateur de fichiers ou le terminal}

\subsection{Flux de traitement}

Le fonctionnement est volontairement séparé en deux étapes, afin de garder une collecte simple et une analyse facile à relancer :

\begin{center}
  \begin{tcolorbox}[colback=codebg, colframe=accent, width=\textwidth, arc=6pt,
    left=12pt, right=12pt, top=10pt, bottom=10pt]
    \color{codetext}\ttfamily\small
    \textcolor{accent}{[Phase 1 — Windows]}\\
    collect.ps1 -TargetFolder "C:\textbackslash{}dossier"\\
    \quad\textcolor{gray}{$\rightarrow$} data/inventory.csv\\
    \quad\textcolor{gray}{$\rightarrow$} data/inventory.json\\[8pt]
    \textcolor{accent}{[Phase 2 — Python]}\\
    python analyze.py --input data/inventory.json\\
    \quad\textcolor{gray}{$\rightarrow$} results/risk\_report.html\\
    \quad\textcolor{gray}{$\rightarrow$} results/risk\_report.json\\
    \quad\textcolor{gray}{$\rightarrow$} results/hashes.txt
  \end{tcolorbox}
\end{center}

% ══════════════════════════════════════════════════════════════════════════════
\section{Script PowerShell — Collecte}

\subsection{Rôle}

Le script \texttt{collect.ps1} parcourt le dossier cible de façon récursive et construit l'inventaire des fichiers trouvés. Il s'appuie uniquement sur des commandes PowerShell natives, ce qui limite les prérequis et facilite l'exécution sur un poste Windows standard.

\subsection{Données collectées}

\begin{table}[h!]
  \centering
  \begin{tabularx}{\textwidth}{lX}
    \toprule
    \textbf{Champ} & \textbf{Description} \\
    \midrule
    \texttt{Name}       & Nom complet du fichier avec extension \\
    \texttt{Extension}  & Extension en minuscules (ex : \texttt{.exe}) \\
    \texttt{SizeBytes}  & Taille en octets \\
    \texttt{CreatedAt}  & Date et heure de création (\texttt{yyyy-MM-dd HH:mm:ss}) \\
    \texttt{ModifiedAt} & Date et heure de dernière modification \\
    \texttt{FullPath}   & Chemin absolu complet sur le disque \\
    \texttt{SHA256}     & Empreinte SHA-256 calculée via \texttt{Get-FileHash} \\
    \bottomrule
  \end{tabularx}
  \caption{Champs collectés par le script PowerShell}
\end{table}

\subsection{Gestion des erreurs}

Certains fichiers peuvent être verrouillés par Windows ou par une application ouverte. Plutôt que d'arrêter toute la collecte, le script isole l'erreur avec un bloc \texttt{Try/Catch}, passe au fichier suivant, puis consigne les éléments non lus dans \texttt{data/skipped\_files.txt}. L'analyste garde ainsi une trace claire des limites du scan.

\subsection{Extrait de code}

\begin{lstlisting}[style=ps1style, caption={Collecte de hash SHA-256 avec gestion d'erreur}]
foreach ($file in $files) {
    try {
        $hash = (Get-FileHash -Path $file.FullName `
                 -Algorithm SHA256 -ErrorAction Stop).Hash
        $entry = [PSCustomObject]@{
            Name       = $file.Name
            Extension  = $file.Extension.ToLower()
            SizeBytes  = $file.Length
            CreatedAt  = $file.CreationTime.ToString("yyyy-MM-dd HH:mm:ss")
            ModifiedAt = $file.LastWriteTime.ToString("yyyy-MM-dd HH:mm:ss")
            FullPath   = $file.FullName
            SHA256     = $hash
        }
        $inventory += $entry
    } catch {
        Write-Warning "Fichier ignore : $($file.FullName)"
        $skipped += [PSCustomObject]@{
            FullPath = $file.FullName
            Reason   = $_.Exception.Message
        }
    }
}
\end{lstlisting}

\subsection{Captures d'écran — Collecte PowerShell}

\screenshot{02\_collect\_ps1\_execution.png}{Exécution de \texttt{collect.ps1} dans PowerShell — affichage de la progression et du nombre de fichiers traités}

\screenshot{03\_inventory\_csv\_apercu.png}{Aperçu du fichier \texttt{data/inventory.csv} ouvert dans Excel ou un éditeur texte — colonnes Name, Extension, SizeBytes, SHA256 visibles}

\screenshot{04\_inventory\_json\_extrait.png}{Extrait de \texttt{data/inventory.json} dans un éditeur — structure JSON avec les champs collectés}

% ══════════════════════════════════════════════════════════════════════════════
\section{Script Python — Analyse et Détection}

\subsection{Architecture de détection}

Le script \texttt{analyze.py} ne cherche pas à remplacer un antivirus. Il combine plutôt plusieurs indices faciles à expliquer : extension, nom, emplacement, dates et réputation du hash. Chaque indice ajoute des points au score final, ce qui permet de prioriser les fichiers à vérifier.

\begin{table}[h!]
  \centering
  \begin{tabularx}{\textwidth}{lXc}
    \toprule
    \textbf{Couche} & \textbf{Règle} & \textbf{Score} \\
    \midrule
    Double extension     & \texttt{fichier.pdf.exe} — extension réelle masquée & +3 \\
    Extension sensible   & \texttt{.exe .ps1 .bat .cmd .vbs .js .scr .lnk}    & +2 \\
    Chemin inhabituel    & \texttt{\textbackslash{}Temp\textbackslash{}}, \texttt{\textbackslash{}AppData\textbackslash{}Roaming\textbackslash{}}, etc. & +2 \\
    Nom suspect          & Mots leurres (invoice, crack, update...) ou imitation de processus système & +2 \\
    Anomalie horodatage  & Modification antérieure à la création (timestomping), création en heure atypique & +2 \\
    \midrule
    Base locale          & Hash présent dans \texttt{known\_malware\_hashes.json}  & +10 \\
    VirusTotal API       & 3+ moteurs antivirus ont détecté le fichier             & +10 \\
    VirusTotal API       & 1--2 moteurs ont détecté le fichier                     & +4  \\
    \bottomrule
  \end{tabularx}
  \caption{Règles de détection et score de risque associé}
\end{table}

\textbf{Niveaux de risque :}
\begin{itemize}[noitemsep]
  \item \textbf{HIGH} — score $\geq$ 3
  \item \textbf{MEDIUM} — score $\geq$ 1
  \item \textbf{LOW} — score = 0
\end{itemize}

\subsection{Rôle du hash SHA-256}

Le SHA-256 sert ici de carte d'identité technique du fichier. Si deux fichiers sont strictement identiques, ils auront la même empreinte ; si le contenu change, même légèrement, l'empreinte change aussi. Dans ce projet, ce hash est utilisé pour deux vérifications :

\begin{enumerate}[itemsep=4pt]
  \item \textbf{Identification par base locale} : le hash est comparé à la liste de hashes connus de la base \texttt{known\_malware\_hashes.json}. Un résultat positif signifie que le fichier est un malware connu, indépendamment de son nom.
  \item \textbf{Corrélation VirusTotal} : le hash est envoyé à l'API VirusTotal v3. La réponse indique combien de moteurs antivirus ont analysé ce fichier exact et l'ont marqué comme malveillant.
\end{enumerate}

\begin{riskbox}[Pourquoi le hash est plus fiable que le nom ?]{accent}
  \color{white}
  Un attaquant peut renommer \texttt{malware.exe} en \texttt{document.pdf}. Le nom devient rassurant, mais le contenu du fichier ne change pas. La comparaison par hash permet donc de dépasser les apparences.
\end{riskbox}

\subsection{Détection des anomalies d'horodatage}

Le timestomping consiste à modifier les dates d'un fichier pour brouiller la chronologie d'une attaque. L'outil ne prétend pas prouver l'intention malveillante à lui seul, mais il signale deux cas qui méritent une vérification :

\begin{itemize}[itemsep=4pt]
  \item \textbf{Modification antérieure à la création} : physiquement impossible dans des conditions normales.
  \item \textbf{Création entre 0h et 5h du matin} : pour les fichiers exécutables, une création nocturne est inhabituelle pour une activité utilisateur légitime.
\end{itemize}

\subsection{Extrait de code}

\begin{lstlisting}[style=pystyle, caption={Détection d'anomalie d'horodatage}]
def check_timestamp_anomaly(entry: dict) -> list[str]:
    hits = []
    created  = datetime.strptime(entry["CreatedAt"],  "%Y-%m-%d %H:%M:%S")
    modified = datetime.strptime(entry["ModifiedAt"], "%Y-%m-%d %H:%M:%S")

    # Timestomping : modification avant creation
    if modified < created:
        hits.append(
            f"Anomalie : modifie le {entry['ModifiedAt']} "
            f"AVANT la creation le {entry['CreatedAt']}"
        )

    # Creation nocturne pour fichier executable
    if entry.get("Extension", "").lower() in SENSITIVE_EXTENSIONS:
        if 0 <= created.hour < 5:
            hits.append(
                f"Cree a {entry['CreatedAt']} "
                f"(heure atypique pour une activite utilisateur)"
            )
    return hits
\end{lstlisting}

\subsection{Captures d'écran — Analyse Python}

\screenshot{05\_analyze\_py\_lancement.png}{Lancement de \texttt{analyze.py} dans le terminal — commande complète avec le paramètre \texttt{--input data/inventory.json}}

\screenshot{06\_analyze\_py\_sortie\_terminal.png}{Sortie terminal de l'analyse en cours : progression fichier par fichier, scores intermédiaires, détections signalées en temps réel}

\screenshot{07\_detection\_double\_extension.png}{Exemple de détection d'une double extension (\texttt{invoice.pdf.exe}) — affichage du score cumulé et des règles déclenchées}

\screenshot{08\_detection\_chemin\_suspect.png}{Exemple de fichier détecté dans un chemin inhabituel (\texttt{Temp} ou \texttt{AppData}) — raison affichée dans la sortie terminal}

% ══════════════════════════════════════════════════════════════════════════════
\section{Intégration VirusTotal}

\subsection{Fonctionnement}

VirusTotal permet de consulter la réputation d'un fichier à partir de son empreinte. Dans ce projet, l'API v3 est interrogée avec le SHA-256, sans téléverser le fichier lui-même.

Cette approche a plusieurs avantages pratiques :
\begin{itemize}[noitemsep]
  \item Le fichier n'est \textbf{pas envoyé} à VirusTotal : seule son empreinte est utilisée.
  \item La réponse est rapide lorsque le hash est déjà connu par la plateforme.
  \item Les fichiers internes ou sensibles restent dans l'environnement de laboratoire.
\end{itemize}

\subsection{Stratégie d'optimisation}

\begin{riskbox}[Gestion du quota API (tier gratuit = 4 req/min)]{accent}
  \color{white}
  \begin{itemize}[noitemsep]
    \item Un hash déjà présent dans la base locale n'est pas redemandé à VirusTotal.
    \item Le script attend 16 secondes entre deux requêtes pour respecter le quota.
    \item En cas d'erreur réseau, l'analyse continue et l'incident est indiqué dans le rapport.
  \end{itemize}
\end{riskbox}

\subsection{Interprétation des résultats}

\begin{table}[h!]
  \centering
  \begin{tabularx}{\textwidth}{lXc}
    \toprule
    \textbf{Résultat VT} & \textbf{Signification} & \textbf{Score} \\
    \midrule
    Hash inconnu (404)   & Fichier jamais soumis à VT — pas nécessairement sain & 0 \\
    0 moteur             & Tous les moteurs concluent à un fichier sain & 0 \\
    1--2 moteurs         & Faux positif possible, à investiguer & +4 \\
    3+ moteurs           & Malware confirmé par consensus & +10 \\
    \bottomrule
  \end{tabularx}
  \caption{Interprétation des résultats VirusTotal}
\end{table}

\subsection{Captures d'écran — VirusTotal}

\screenshot{09\_virustotal\_requete\_terminal.png}{Sortie terminal lors de l'appel à l'API VirusTotal — affichage du hash envoyé, du temps d'attente et de la réponse reçue (nombre de moteurs ayant détecté le fichier)}

\screenshot{10\_virustotal\_interface\_web.png}{Interface web de VirusTotal montrant le résultat pour un hash testé — onglet \textit{Detection} avec le score moteurs/total (optionnel, pour illustration)}

% ══════════════════════════════════════════════════════════════════════════════
\section{Résultats et rapports générés}

\subsection{Fichiers de sortie}

À la fin de l'analyse, les résultats sont écrits dans \texttt{results/}. Les formats ont été choisis pour couvrir les usages courants : lecture rapide, archivage et réutilisation par d'autres outils.

\begin{itemize}[itemsep=6pt]
  \item \textbf{\texttt{risk\_report.html}} — Rapport interactif avec interface sombre, cartes par fichier, filtres par niveau de risque, barre de score visuelle, bouton de copie du hash, et liens cliquables vers VirusTotal.
  \item \textbf{\texttt{risk\_report.txt}} — Rapport texte brut lisible dans un terminal ou un éditeur.
  \item \textbf{\texttt{risk\_report.json}} — Rapport structuré complet, exploitable par des outils tiers.
  \item \textbf{\texttt{hashes.txt}} — Liste des empreintes SHA-256 au format \texttt{sha256sum}.
\end{itemize}

\subsection{Exemple de sortie terminal}

\begin{lstlisting}[style=ps1style, caption={Exemple de sortie lors de l'analyse}]
[1/3] invoice.pdf.exe
    Path    : C:\Users\Jihene\Desktop\virus\invoice.pdf.exe
    Size    : 12 bytes
    Created : 2024-01-15 02:34:00
    Modified: 2024-01-10 09:00:00
    SHA-256 : 3a9f2c1b7e...
    Checks  :
      [!] +3  Double extension : invoice.pdf.exe
      [!] +2  Extension executable : .exe
      [ ]     Chemin normal
      [!] +2  Mot leurre : 'invoice'
      [!] +2  Anomalie horodatage : modifie avant creation
      [ ]     Hash absent de la base locale
      [~]     VirusTotal... 5/72 moteurs - CONFIRME MALWARE  +10
    Result  : HIGH (score: 19)
\end{lstlisting}

\subsection{Captures d'écran — Rapports générés}

\screenshot{11\_rapport\_html\_vue\_globale.png}{Vue globale du rapport \texttt{risk\_report.html} ouvert dans le navigateur — en-tête du tableau de bord avec les compteurs HIGH / MEDIUM / LOW}

\screenshot{12\_rapport\_html\_carte\_high.png}{Carte d'un fichier classé HIGH dans le rapport HTML — score visible, règles déclenchées listées, lien VirusTotal cliquable}

\screenshot{13\_rapport\_html\_filtre\_risque.png}{Démonstration du filtre par niveau de risque dans l'interface HTML — filtrage sur HIGH uniquement}

\screenshot{14\_rapport\_txt\_terminal.png}{Affichage du fichier \texttt{risk\_report.txt} dans le terminal (\texttt{cat results/risk\_report.txt}) — résumé lisible par un analyste}

\screenshot{15\_hashes\_txt.png}{Contenu du fichier \texttt{hashes.txt} — liste des empreintes SHA-256 au format \texttt{sha256sum}, une ligne par fichier analysé}

\screenshot{16\_results\_dossier.png}{Dossier \texttt{results/} après l'exécution complète — les quatre fichiers de sortie présents (\texttt{.html}, \texttt{.json}, \texttt{.txt}, \texttt{hashes.txt})}

% ══════════════════════════════════════════════════════════════════════════════
\section{Conclusion et recommandations de sécurité}

\subsection{Bilan}

Ce projet aboutit à un outil simple, mais déjà utile pour une première investigation défensive. PowerShell prend en charge la collecte sur Windows, tandis que Python apporte la logique d'analyse, de scoring et de génération des rapports.

L'intérêt principal vient de la combinaison des indices. Une extension suspecte seule ne suffit pas toujours à conclure ; en revanche, une double extension, un chemin inhabituel, une anomalie de date et un hash détecté donnent une priorité claire à l'analyste.

L'intégration VirusTotal apporte une dimension de réputation externe, complémentaire à la base locale. Le pipeline complet — collecte PowerShell, analyse Python, rapport HTML — peut s'exécuter en quelques minutes sur un dossier de taille raisonnable, ce qui le rend exploitable dans un contexte de tri rapide.

\screenshot{17\_bilan\_execution\_complete.png}{Résumé final affiché dans le terminal à la fin de l'analyse — nombre total de fichiers traités, répartition HIGH / MEDIUM / LOW, temps d'exécution}

\subsection{Recommandations pour l'analyste}

\begin{riskbox}[Recommandations pour l'analyste]{accent}
  \color{white}
  \begin{enumerate}[itemsep=6pt]
    \item Ne jamais ouvrir un fichier signalé sur une machine de production ; utiliser une sandbox ou une machine virtuelle.
    \item Traiter les fichiers \textbf{HIGH} en premier, puis vérifier les cas \textbf{MEDIUM} selon le contexte.
    \item Conserver les hashes confirmés malveillants pour enrichir progressivement la base locale \texttt{known\_malware\_hashes.json}.
    \item Examiner les anomalies d'horodatage avec prudence : elles sont utiles, mais doivent être recoupées avec d'autres indices.
    \item Porter une attention particulière aux exécutables placés dans \texttt{Temp}, \texttt{AppData} ou d'autres chemins rarement utilisés par l'utilisateur.
    \item Documenter chaque investigation : conserver les fichiers \texttt{inventory.json} et \texttt{risk\_report.json} comme preuve d'audit.
  \end{enumerate}
\end{riskbox}

\subsection{Recommandations pour renforcer l'outil}

\begin{riskbox}[Recommandations pour renforcer l'outil]{accent}
  \color{white}
  \begin{enumerate}[itemsep=6pt]
    \item Ajouter une phase d'analyse dynamique en sandbox pour observer le comportement réel des fichiers suspects avant toute décision.
    \item Intégrer des règles YARA afin de détecter des familles ou motifs connus sans dépendre d'une connexion réseau.
    \item Prévoir des scans réguliers et automatisés sur les dossiers sensibles (\texttt{Downloads}, \texttt{Temp}, \texttt{AppData}).
    \item Exporter les alertes \textbf{HIGH} vers un SIEM (Splunk, Elastic) pour faciliter le suivi et la corrélation avec d'autres événements.
    \item Enrichir la base locale \texttt{known\_malware\_hashes.json} avec des flux de réputation publics comme MalwareBazaar ou CIRCL Hashlookup.
    \item Ajouter une signature ou un hash du rapport lui-même pour garantir son intégrité lors d'un audit ultérieur.
  \end{enumerate}
\end{riskbox}

\subsection{Limites actuelles}

\begin{itemize}[itemsep=4pt]
  \item L'analyse reste statique : un malware obfusqué ou chiffré peut passer inaperçu si son hash n'est pas référencé.
  \item Le quota gratuit de VirusTotal limite le nombre de requêtes par minute ; un grand dossier ralentit l'analyse.
  \item Certains scripts légitimes peuvent être signalés à tort ; la décision finale revient toujours à l'analyste.
  \item Le script PowerShell ne récupère pas les flux de données alternatifs (NTFS ADS), qui peuvent dissimuler du contenu malveillant.
\end{itemize}

% ── Bibliographie ─────────────────────────────────────────────────────────────
\newpage
\section*{Références}
\addcontentsline{toc}{section}{Références}
\small
\sloppy

\begin{itemize}[itemsep=6pt]
  \item Documentation de l'API VirusTotal v3 — \url{https://developers.virustotal.com/reference}
  \item Fichier de test antivirus EICAR — \url{https://www.eicar.org}
  \item MalwareBazaar — \url{https://bazaar.abuse.ch}
  \item CIRCL Hashlookup — \url{https://hashlookup.circl.lu}
  \item MITRE ATT\&CK — Timestomping (T1070.006) — \url{https://attack.mitre.org/techniques/T1070/006/}
  \item Documentation PowerShell, \texttt{Get-FileHash} — \url{https://learn.microsoft.com/powershell/module/microsoft.powershell.utility/get-filehash}
  \item Documentation Python, \texttt{pathlib} — \url{https://docs.python.org/3/library/pathlib.html}
\end{itemize}
\normalsize
\fussy

\end{document}
